\section{Fácies, associações de fácies e ambientes deposicionais}

\begin{definicao}
Uma \textbf{fácies sedimentar} é um corpo de sedimento ou rocha caracterizado por um
conjunto de propriedades observáveis que o distingue dos corpos adjacentes. Essas
propriedades podem incluir litologia, granulometria, estruturas sedimentares, fósseis,
cor, geometria e contatos.
\end{definicao}

A descrição de uma fácies deve preceder sua interpretação. Por exemplo, ``arenito fino
com laminação cruzada e bioturbação esparsa'' é uma descrição; ``depósito de canal de
maré'' é uma interpretação que precisa ser sustentada pelo conjunto de evidências.

\subsection{Códigos de litofácies}

Os códigos de litofácies combinam uma letra maiúscula que indica a classe granulométrica
com letras minúsculas que representam estruturas ou características. Em esquemas
derivados de Miall, por exemplo, \emph{G} indica cascalho, \emph{S} indica areia e
\emph{F} indica sedimentos finos. O significado completo de cada código deve ser
declarado no trabalho, porque há adaptações entre autores e ambientes
\parencite{miall1978,nichols2009}.

\subsection{Associação de fácies}

Uma fácies isolada raramente define um ambiente deposicional de maneira inequívoca. Uma
\textbf{associação de fácies} reúne fácies que ocorrem repetidamente juntas e que podem
ser relacionadas por processos sedimentares compatíveis. A interpretação ganha força
quando considera:

\begin{itemize}
  \item a sucessão vertical das fácies;
  \item as relações laterais;
  \item a geometria dos corpos;
  \item as superfícies erosivas ou de reativação;
  \item a direção das paleocorrentes;
  \item os fósseis e icnofósseis;
  \item o contexto estratigráfico regional.
\end{itemize}

\subsection{Sucessão de fácies}

Uma \textbf{sucessão de fácies} é a ordem vertical ou lateral em que diferentes fácies
aparecem. Tendências granodecrescentes podem registrar diminuição da energia do fluxo,
como em certos preenchimentos de canal. Tendências granocrescentes podem acompanhar a
progradação de sistemas deposicionais. Essas interpretações, entretanto, dependem da
geometria, das superfícies limitantes e do ambiente analisado.

\subsection{Ambientes continentais, transicionais e marinhos}

Em escala ampla, os ambientes deposicionais podem ser agrupados em:

\begin{itemize}
  \item \textbf{continentais}: fluviais, lacustres, eólicos, glaciais e leques aluviais;
  \item \textbf{transicionais}: deltas, estuários, planícies de maré e sistemas costeiros;
  \item \textbf{marinhos}: plataforma, talude, bacia profunda, recifes e outros sistemas
        oceânicos.
\end{itemize}

Estratificação espinha de peixe, superfícies de reativação e alternância de laminações de
areia e lama podem indicar reversão periódica de correntes, comum em ambientes de maré.
Camadas com estratificação cruzada hummocky são associadas a fluxos combinados gerados
por tempestades em ambientes marinhos rasos. Nenhuma dessas feições deve ser usada
isoladamente como diagnóstico definitivo \parencite{harms1975,nichols2009}.
